\documentclass[11pt,a4paper]{article}
\usepackage[utf8]{inputenc}
\usepackage{amsmath,amssymb,amsthm}
\usepackage{listings}
\usepackage[margin=2.5cm]{geometry}
\usepackage{xcolor}

\lstset{
  language=C++,
  basicstyle=\ttfamily\small,
  keywordstyle=\color{blue}\bfseries,
  commentstyle=\color{green!50!black},
  stringstyle=\color{red!70!black},
  numbers=left,
  numberstyle=\tiny\color{gray},
  breaklines=true,
  frame=single,
  tabsize=4,
  showstringspaces=false
}

\title{IOI 1993 -- Day 1, Task 1: The Castle}
\author{Solution}
\date{}

\begin{document}
\maketitle

\section{Problem Statement}

A castle is represented as an $M \times N$ grid (where the input gives $M$
columns and $N$ rows). Each cell's walls are encoded as a 4-bit integer:
\begin{itemize}
  \item Bit 0 (value 1): wall on the West.
  \item Bit 1 (value 2): wall on the North.
  \item Bit 2 (value 4): wall on the East.
  \item Bit 3 (value 8): wall on the South.
\end{itemize}

A \emph{room} is a maximal connected region of cells not separated by walls.
Determine:
\begin{enumerate}
  \item The number of rooms.
  \item The area of the largest room.
  \item The largest room obtainable by removing a single wall, and which
    wall to remove.
\end{enumerate}

\section{Solution}

\subsection{Phase 1: Room Identification (Flood Fill)}

Two adjacent cells $(r, c)$ and $(r', c')$ are connected if no wall
separates them. The wall bitmask makes adjacency checks direct:

\begin{center}
\begin{tabular}{lll}
Direction & Neighbor & No wall if \\
\hline
West  & $(r, c-1)$ & bit 0 is clear \\
North & $(r-1, c)$ & bit 1 is clear \\
East  & $(r, c+1)$ & bit 2 is clear \\
South & $(r+1, c)$ & bit 3 is clear \\
\end{tabular}
\end{center}

BFS from each unvisited cell assigns room IDs and computes room sizes.

\subsection{Phase 2: Best Wall Removal}

For every internal wall between two cells in different rooms, the combined
room size is $\text{size}[\text{room}_1] + \text{size}[\text{room}_2]$.
Track the maximum.

The tie-breaking rule (per the USACO specification) requires iterating
column-by-column (west to east), bottom to top within each column, and
checking North walls before East walls.

\section{Complexity Analysis}

\begin{itemize}
  \item \textbf{Time:} $O(MN)$ for flood fill and $O(MN)$ for wall
    enumeration.
  \item \textbf{Space:} $O(MN)$.
\end{itemize}

\section{Implementation}

\begin{lstlisting}
#include <iostream>
#include <vector>
#include <queue>
#include <algorithm>
#include <cstring>
using namespace std;

int M, N; // M = columns, N = rows
int grid[55][55];
int roomId[55][55];
int roomSize[2505];
int numRooms;

// Directions: W, N, E, S
const int dr[] = {0, -1, 0, 1};
const int dc[] = {-1, 0, 1, 0};
const int wallBit[] = {1, 2, 4, 8};

void bfs(int sr, int sc, int id) {
    queue<pair<int,int>> q;
    q.push({sr, sc});
    roomId[sr][sc] = id;
    int sz = 0;

    while (!q.empty()) {
        auto [r, c] = q.front(); q.pop();
        sz++;
        for (int d = 0; d < 4; d++) {
            if (grid[r][c] & wallBit[d]) continue;
            int nr = r + dr[d], nc = c + dc[d];
            if (nr < 0 || nr >= N || nc < 0 || nc >= M) continue;
            if (roomId[nr][nc] != -1) continue;
            roomId[nr][nc] = id;
            q.push({nr, nc});
        }
    }
    roomSize[id] = sz;
}

int main() {
    cin >> M >> N;

    for (int r = 0; r < N; r++)
        for (int c = 0; c < M; c++)
            cin >> grid[r][c];

    // Phase 1: find rooms
    memset(roomId, -1, sizeof(roomId));
    numRooms = 0;

    for (int r = 0; r < N; r++)
        for (int c = 0; c < M; c++)
            if (roomId[r][c] == -1)
                bfs(r, c, numRooms++);

    int maxRoom = 0;
    for (int i = 0; i < numRooms; i++)
        maxRoom = max(maxRoom, roomSize[i]);

    cout << numRooms << endl;
    cout << maxRoom << endl;

    // Phase 2: find best wall to remove
    // Iterate: column by column (west to east),
    //          bottom to top within each column,
    //          check N wall before E wall.
    int bestSize = 0;
    int bestR = -1, bestC = -1;
    char bestDir = ' ';

    for (int c = 0; c < M; c++) {
        for (int r = N - 1; r >= 0; r--) {
            // Check North wall
            if (r > 0 && (grid[r][c] & 2)) {
                int id1 = roomId[r][c];
                int id2 = roomId[r-1][c];
                if (id1 != id2) {
                    int combined = roomSize[id1] + roomSize[id2];
                    if (combined > bestSize) {
                        bestSize = combined;
                        bestR = r; bestC = c; bestDir = 'N';
                    }
                }
            }
            // Check East wall
            if (c < M - 1 && (grid[r][c] & 4)) {
                int id1 = roomId[r][c];
                int id2 = roomId[r][c+1];
                if (id1 != id2) {
                    int combined = roomSize[id1] + roomSize[id2];
                    if (combined > bestSize) {
                        bestSize = combined;
                        bestR = r; bestC = c; bestDir = 'E';
                    }
                }
            }
        }
    }

    cout << bestSize << endl;
    cout << bestR + 1 << " " << bestC + 1 << " " << bestDir << endl;

    return 0;
}
\end{lstlisting}

\section{Example}

\begin{verbatim}
Input:
7 4
11 6 11 6 3 10 6
7  9  6 13 5 15 5
1 10 12  7 13  7 5
13 11 10  8 10 12 13

Output:
5        (number of rooms)
9        (largest room)
16       (largest after removing one wall)
4 1 E    (row 4, column 1, East wall)
\end{verbatim}

\section{Notes}

\begin{itemize}
  \item The wall bitmask encoding is compact: a single integer per cell
    fully describes all four walls.
  \item Consistency is guaranteed by the input: if cell $(r,c)$ has a wall
    on the East, then cell $(r, c+1)$ has a wall on the West.
  \item This problem also appears in the USACO Training Pages and is a
    classic application of flood fill with bitmask wall encoding.
\end{itemize}

\end{document}
